% FILE: 04_implementation_surface.tex
% Project: schemas
% Thesis Role: data-schema-definitions-for-ma-and-receipt-surfaces

\section{Implementation Surface}
\label{sec:schemas-implementation}

\subsection{Source Files}
\label{sec:schemas-implementation-source}

The implementation surface of the schemas project is deliberately minimal. Two JSON Schema files
constitute the entire normative surface:

\begin{description}
  \item[\texttt{schemas/slide\_verification\_ledger.schema.json}] 65 lines. Defines the
    \texttt{SlideVerificationLedger} type using JSON Schema draft-07. Uses
    \texttt{patternProperties} with a UUIDv4 regex key pattern to admit only UUID-identified
    slide receipts. Eight required fields per receipt entry. Strict enforcement via
    \texttt{additionalProperties: false}.
  \item[\texttt{standards/otel\_trace\_integration\_schema.json}] 75 lines. Defines the
    \texttt{OTelTraceBlake3IntegrationSchema} type. Requires W3C TraceID, BLAKE3 event-chain root,
    and a spans array. Each span element requires W3C SpanID, parent pointer, activity label,
    microsecond timestamps, cumulative instruction count, and a per-span BLAKE3 receipt.
\end{description}

Five production receipt files in \texttt{receipts/} implement the schema contract:

\begin{description}
  \item[\texttt{rec\_ebitda\_rework\_001.json}] Slide 1 receipt. Asserts \$1,250,000 annual EBITDA
    improvement. Engine: \texttt{wasm4pm}. Query URI:
    \texttt{file:///\ldots/queries/ebitda\_rework\_reduction.wasm}. Fitness: 0.982. Precision: 0.945.
    BLAKE3 hashes: \texttt{ee9ab723\ldots} (log), \texttt{81f7dca2\ldots} (model).
  \item[\texttt{rec\_wc\_ar\_002.json}] Slide 2 receipt. Asserts \$1,369,863 working capital
    release via AR velocity acceleration (42.5-day to 32.5-day cycle time reduction on
    \$50M annual credit sales). Fitness: 0.978. Precision: 0.923.
  \item[\texttt{rec\_risk\_sla\_003.json}] Slide 3 receipt. Asserts SLA penalty liability capped
    at \$450,000 with late delivery rates below 2.5\%.
  \item[\texttt{rec\_risk\_compliance\_004.json}] Slide 4 receipt. Asserts \$0 compliance leakage
    with zero SoD or SOX/GDPR/AML violations in procurement workflows.
  \item[\texttt{rec\_residual\_standard\_005.json}] Slide 5 receipt. Asserts 97.5\% conformance
    standardization with residual weight $W_R \le 0.025$ and residual entropy $H_R = 0.85$.
\end{description}

All five receipts carry \texttt{validator\_signature} values prefixed \texttt{ed25519\_sig\_} with
UUIDs and 64-byte hex payloads, conforming to the 128-hex pattern. All five reference
\texttt{\$schema: https://json-schema.org/draft/2020-12/schema} (using the draft-2020-12 URI
despite the schema file itself targeting draft-07---a minor inconsistency in the sample receipts
that is noted but does not affect structural conformance of the fields).

\subsection{Commands}
\label{sec:schemas-implementation-commands}

No build system, CLI tool, or automated validation command is attached to the schemas project at
time of assessment. Schema validation must be performed manually or via an external tool such as
\texttt{ajv} (Node.js) or \texttt{jsonschema} (Python). This is the primary gap in the
implementation surface: the contract exists but is not enforced programmatically at commit time.

The M\&A deck rendering authority assessment (2026-06-01) defines a seven-step rendering pipeline
that includes a validation step (Step~6) which verifies receipt structure against the schema.
However, this validation step is defined as a protocol specification, not as an executed command
with a recorded output. The command that would implement it does not exist as a committed artifact.

\AUTHORINTERPRETATION{} The absence of an automated validation command is the single most
consequential structural gap in the schemas project. It means the schema is a specification, not
an enforced law.

\subsection{Data and Ontology Surfaces}
\label{sec:schemas-implementation-data}

The schemas project intersects the ontology surface of the research program at the
\texttt{query\_definition.engine} field. The enum values \texttt{wasm4pm} and \texttt{pm4py}
name the two execution engines recognized by the research program. This naming is consistent with
the wasm4pm authority maps in \texttt{sources/wasm4pm/} and the pm4py oracle surface in
\texttt{sources/pm4py/}.

The BLAKE3 hash fields intersect the cryptographic receipt doctrine defined in
\texttt{doctrine/RECEIPT\_DOCTRINE.md}. That document establishes that a receipt must be a
typed, witnessed, bound evidence artifact. The schema operationalizes ``bound'' by requiring
that the receipt carry content addresses (BLAKE3 hashes) of both the event log and the process
model it references.

The JCS canonical serialization requirement (RFC~8785) for the Ed25519 signature computation is
specified in the slide-to-receipt map (\texttt{ma/slide-to-receipt-map.md}) but not in the schema
itself. The schema pattern for \texttt{validator\_signature} (\verb|^[0-9a-fA-F]{128}$|) admits
any 128-hex string, without requiring that the string represent a valid Ed25519 signature over
a JCS-canonical payload. The binding between the pattern and the signing protocol is therefore
extra-schematic---a documentation gap.

\subsection{Templates and Rendered Artifacts}
\label{sec:schemas-implementation-templates}

The schema is consumed by two Tera2 templates:

\begin{description}
  \item[\texttt{ggen/templates/ma-deck.tera}] (202 lines) The PowerPoint rendering template.
    References receipt hashes (first 16 characters) in slide-level claim detail slides. Computes
    board admissibility as \texttt{status = "ADMISSIBLE"} if \texttt{verdictFitness} $\ge 0.95$
    and \texttt{verdictPrecision} $\ge 0.90$; else \texttt{"REJECTED"}.
  \item[\texttt{ggen/templates/ma-diligence.tera}] (283 lines) The Excel rendering template.
    Exposes receipt hashes in the Synergy\_Claims and Replay\_Traces sheets, enabling auditors to
    cross-reference the Excel workbook against the receipt files in the VDR.
\end{description}

Both templates consume claims data as a JSON array whose \texttt{receiptHash} and \texttt{receipt}
fields reference the ledger. The templates do not validate schema conformance---that is the
responsibility of Step~6 in the rendering pipeline. The templates assume that any receipt hash
they receive has already been validated against the schema.
